\documentclass[aspectratio=169,11pt]{beamer}
\usepackage[indonesian]{babel}
\usepackage{booktabs,tikz,listings,amsmath,hyperref,multicol}
\usetheme{metropolis}
\usebeamercolor{normal text}\colorlet{codeink}{normal text.fg}
\usebeamercolor{alerted text}\colorlet{codeaccent}{alerted text.fg!60!codeink}
\lstset{language={[LaTeX]TeX},basicstyle=\ttfamily\small\color{codeink},
  texcsstyle=*\bfseries\color{codeaccent},columns=fullflexible,keepspaces=true,
  showstringspaces=false,backgroundcolor=\color{codeink!6},frame=l,
  rulecolor=\color{codeaccent},xleftmargin=5pt,aboveskip=4pt,belowskip=4pt}
\lstdefinestyle{bib}{basicstyle=\ttfamily\scriptsize,breaklines=true}
\lstdefinestyle{terminal}{language=bash,basicstyle=\ttfamily\small\color{codeink},breaklines=true}
\newcommand{\tugas}[1]{\par\medskip\alert{\textbf{Coba sendiri}}\quad #1}
\title{Menulis dokumen akademik dengan \LaTeX}
\subtitle{Hari 3 --- Tata letak, kustomisasi, dan finalisasi}
\author{Pelatihan Universitas Terbuka dan Departemen Matematika FMIPA UGM}
\date{Yogyakarta, 19 November 2026}
\begin{document}
\begin{frame}[plain]\usebeamerfont{title}\inserttitle\par\usebeamerfont{subtitle}\insertsubtitle\par\medskip\usebeamertemplate{title separator}\par\medskip\usebeamerfont{author}\insertauthor\par\smallskip\usebeamerfont{date}\insertdate\end{frame}
\begin{frame}{Tujuan belajar Hari 3}
\begin{enumerate}
\item Mengendalikan tata letak dokumen panjang secara konsisten.
\item Membaca dan memodifikasi template, kelas, paket, perintah, dan lingkungan.
\item Menambahkan bibliografi, indeks, glossary, kode, bahasa, dan materi khusus.
\item Memeriksa sumber, layout, referensi, dan hasil akhir sebelum diserahkan.
\end{enumerate}
\end{frame}
\begin{frame}{Peta belajar Hari 3}
\tableofcontents
\end{frame}
\section{Tata letak dokumen}
\begin{frame}{Tata letak dokumen panjang}
Tata letak yang baik membuat pembaca dapat mengikuti struktur tanpa menebak.
\medskip
\begin{itemize}
\item ukuran kertas dan margin mengikuti kebutuhan institusi;
\item header, footer, dan nomor halaman memberi orientasi;
\item judul, paragraf, daftar, tabel, dan gambar memakai gaya yang berulang;
\item semua keputusan format disimpan di prambule atau berkas gaya.
\end{itemize}
\end{frame}
\begin{frame}[fragile]{Ukuran kertas dan margin}
\begin{lstlisting}
\documentclass[a4paper,12pt]{book}
\usepackage[a4paper,margin=3cm]{geometry}
% atau margin berbeda untuk sisi dalam dan luar
\usepackage[inner=3.5cm,outer=2.5cm,top=3cm,bottom=3cm]{geometry}
\end{lstlisting}
Gunakan satu sumber kebenaran untuk ukuran halaman. Periksa juga area cetak,
ruang jilid, dan persyaratan template institusi.
\tugas{Bandingkan PDF dengan margin 2 cm dan 3 cm.}
\end{frame}
\begin{frame}[fragile]{Single-sided dan double-sided}
\begin{lstlisting}
\documentclass[oneside]{book}  % halaman semua satu sisi
\documentclass[twoside,openright]{book}
\cleardoublepage                % bab baru pada halaman kanan
\end{lstlisting}
Pada mode dua sisi, halaman ganjil dan genap dapat memiliki header berbeda.
Sisakan ruang jilid pada margin dalam dan uji halaman kosong yang sengaja muncul.
\end{frame}
\begin{frame}[fragile]{Header, footer, dan nomor halaman}
\begin{lstlisting}
\usepackage{fancyhdr}
\pagestyle{fancy}
\fancyhead[LE]{\nouppercase{\leftmark}}
\fancyhead[RO]{\nouppercase{\rightmark}}
\fancyfoot[C]{\thepage}
\fancypagestyle{plain}{\fancyhf{}\fancyfoot[C]{\thepage}}
\end{lstlisting}
Jangan menimpa gaya halaman judul tanpa sengaja. Untuk dokumen akademik,
cek apakah nomor halaman awal memakai angka Romawi dan isi memakai angka Arab.
\end{frame}
\begin{frame}[fragile]{Paragraf dan jarak vertikal}
\begin{lstlisting}
\usepackage{parskip}
\setlength{\parindent}{0pt}
\setlength{\parskip}{6pt plus 2pt minus 1pt}
\end{lstlisting}
Pilih satu pola: indentasi paragraf atau jarak antarparagraf. Hindari menambah
\lstinline|\\| berulang kali untuk mengatur jarak karena hasilnya rapuh.
\tugas{Buat dua paragraf dan bandingkan \texttt{parskip} dengan indentasi.}
\end{frame}
\begin{frame}[fragile]{Multiple columns}
\begin{lstlisting}
\usepackage{multicol}
\begin{multicols}{2}
Teks akan mengalir dari kolom kiri ke kolom kanan.
\end{multicols}
\end{lstlisting}
Kolom cocok untuk daftar istilah atau ringkasan. Hindari kolom sempit untuk
paragraf panjang dan pastikan judul, tabel, serta gambar tidak terpotong.
\end{frame}
\begin{frame}[fragile]{Counters dan penomoran}
\begin{lstlisting}
\newcounter{contoh}
\setcounter{contoh}{0}
\newcommand{\contoh}{\stepcounter{contoh}Contoh \thecontoh}
\renewcommand{\thefigure}{\thesection.\arabic{figure}}
\end{lstlisting}
Gunakan counter untuk sesuatu yang benar-benar perlu bernomor. Biarkan LaTeX
mengatur nomor tabel, gambar, persamaan, dan bagian secara otomatis.
\end{frame}
\begin{frame}[fragile]{Daftar isi dan daftar gambar}
\begin{lstlisting}
\tableofcontents
\listoffigures
\listoftables
\section[Singkat di daftar]{Judul bagian yang panjang di halaman}
\caption[Judul singkat]{Keterangan lengkap gambar}
\end{lstlisting}
Kompilasi beberapa kali atau gunakan \texttt{latexmk} agar berkas bantu diperbarui.
Judul pendek menjaga daftar tetap terbaca.
\end{frame}
\begin{frame}{Konsistensi gaya}
Tetapkan keputusan sebelum menulis banyak halaman:
\begin{itemize}
\item satuan margin, ukuran heading, dan jarak antarbagian;
\item bentuk caption, label tabel, dan penulisan istilah;
\item aturan untuk kutipan, daftar, URL, dan nomor persamaan;
\item warna dan jenis huruf yang dapat dibaca saat dicetak hitam putih.
\end{itemize}
\end{frame}
\begin{frame}[fragile]{Pengelolaan proyek besar}
\begin{lstlisting}
main.tex
chapters/pendahuluan.tex
chapters/metode.tex
figures/alat.pdf
references.bib
\input{chapters/pendahuluan}
\include{chapters/metode}
\end{lstlisting}
Simpan satu sumber utama untuk urutan dokumen. Pisahkan isi, gambar,
referensi, dan konfigurasi agar perubahan mudah ditinjau.
\end{frame}
\section{Template dan kustomisasi}
\begin{frame}{Template, kelas, dan paket}
\begin{description}
\item[Template] susunan awal dan gaya yang siap digunakan.
\item[Kelas \texttt{.cls}] aturan dasar jenis dokumen, misalnya artikel atau buku.
\item[Paket \texttt{.sty}] fitur tambahan yang dapat dimuat sesuai kebutuhan.
\item[Prambule] tempat untuk konfigurasi yang berlaku di seluruh dokumen.
\end{description}
\end{frame}
\begin{frame}[fragile]{Memodifikasi template dengan aman}
\begin{lstlisting}
% Jangan mengubah file .cls asli.
\documentclass{template-institusi}
\usepackage{microtype}
\usepackage{csquotes}
\newcommand{\namauniversitas}{Universitas Terbuka}
\end{lstlisting}
Buat berkas konfigurasi sendiri atau gunakan opsi kelas. Catat perubahan
yang diperlukan agar template dapat dipakai ulang oleh peserta lain.
\end{frame}
\begin{frame}[fragile]{Mendefinisikan commands}
\begin{lstlisting}
\newcommand{\hasilutama}[1]{\textbf{#1}}
\newcommand{\R}{\mathbb{R}}
\DeclareMathOperator{\argmin}{arg\,min}
\hasilutama{Kesimpulan utama} berada di $\R^n$.
\end{lstlisting}
Perintah mengurangi pengulangan dan menjaga istilah tetap seragam.
Gunakan nama yang jelas dan definisikan setiap perintah hanya sekali.
\end{frame}
\begin{frame}[fragile]{Membuat environments}
\begin{lstlisting}
\newenvironment{catatan}[1][Catatan]
  {\begin{quote}\textbf{#1.}\ }
  {\end{quote}}
\begin{catatan}[Perhatian]
Periksa sumber data sebelum menyimpulkan.
\end{catatan}
\end{lstlisting}
Environment cocok untuk pola isi yang berulang. Pilih struktur yang
semantis, bukan sekadar cara cepat mengubah jarak.
\end{frame}
\begin{frame}{Gaya institusi dan proyek multi-file}
Gaya institusi sebaiknya memusatkan identitas visual dan aturan format.
\medskip
\begin{itemize}
\item buat paket atau berkas konfigurasi untuk warna, logo, dan heading;
\item simpan bab dalam berkas terpisah dengan \texttt{\\input} atau \texttt{\\include};
\item gunakan Git atau riwayat versi untuk meninjau perubahan;
\item dokumentasikan dependensi dan cara kompilasi.
\end{itemize}
\end{frame}
\begin{frame}[fragile]{Hyperlinks dan format akademik}
\begin{lstlisting}
\usepackage[hidelinks]{hyperref}
\hypersetup{pdftitle={Laporan penelitian},pdfauthor={Nama}}
\href{https://www.example.org}{Situs resmi}
\url{https://www.example.org}
\end{lstlisting}
Gunakan label dan rujukan internal untuk navigasi. Uji tautan setelah
memindahkan bagian atau mengganti nama berkas.
\end{frame}
\section{Fitur penerbitan dan materi khusus}
\begin{frame}[fragile]{Bibliography dengan BibTeX atau biblatex}
\begin{lstlisting}
\usepackage[backend=biber,style=authoryear]{biblatex}
\addbibresource{references.bib}
Menurut \textcite{lamport1994}, struktur sumber penting.
\printbibliography
\end{lstlisting}
Pilih satu alur kerja sesuai template. Jangan mencampur perintah BibTeX
lama dan biblatex tanpa alasan yang jelas.
\end{frame}
\begin{frame}[fragile]{Berkas BibTeX dan gaya sitasi}
\begin{lstlisting}[style=bib]
@book{lamport1994,
  author    = {Lamport, Leslie},
  title     = {LaTeX: A Document Preparation System},
  year      = {1994},
  publisher = {Addison-Wesley}
}
\end{lstlisting}
Gaya sitasi ditentukan oleh disiplin dan panduan penerbit. Periksa nama,
tahun, judul, DOI, dan URL sebelum menyerahkan dokumen.
\end{frame}
\begin{frame}{Indeks, glossary, dan nomenclature}
\begin{itemize}
\item indeks membantu pembaca menemukan istilah di banyak halaman;
\item glossary menjelaskan istilah teknis dan singkatannya;
\item nomenclature mengumpulkan simbol beserta definisinya.
\end{itemize}
Ketiganya memerlukan langkah kompilasi tambahan. Pastikan istilah yang
ditampilkan memang dipakai dan urutannya konsisten.
\end{frame}
\begin{frame}[fragile]{Kode dengan minted}
\begin{lstlisting}
\usepackage{minted}
\begin{minted}[fontsize=\small,linenos]{python}
for value in data:
    print(value)
\end{minted}
\end{lstlisting}
\texttt{minted} memakai Pygments dan biasanya memerlukan opsi kompilasi
\texttt{-shell-escape}. Gunakan hanya pada lingkungan yang dipercaya.
\end{frame}
\begin{frame}[fragile]{Kutipan dan tanda kutip}
\begin{lstlisting}
\usepackage{csquotes}
\enquote{Kutipan pendek mengikuti aturan bahasa.}
\begin{displayquote}
Kutipan panjang ditempatkan sebagai blok tersendiri.
\end{displayquote}
\end{lstlisting}
Jangan mengetik tanda kutip lurus sebagai pengganti tanda kutip tipografis.
Panjang kutipan dan aturan indentasi mengikuti gaya publikasi.
\end{frame}
\begin{frame}[fragile]{Bahasa Indonesia dan bahasa internasional}
\begin{lstlisting}
\usepackage[indonesian,english]{babel}
\selectlanguage{indonesian}
\foreignlanguage{english}{Abstract}
\end{lstlisting}
Periksa pemenggalan kata, nama bulan, judul otomatis, tanda baca, dan
aturan huruf kapital. Pilih font yang mendukung semua karakter yang dipakai.
\end{frame}
\begin{frame}[fragile]{Theorems dan proofs}
\begin{lstlisting}
\usepackage{amsthm}
\newtheorem{teorema}{Teorema}[section]
\begin{teorema}
Jumlah dua bilangan genap adalah genap.
\end{teorema}
\begin{proof}
Tuliskan alasan matematis secara lengkap.
\end{proof}
\end{lstlisting}
Gunakan lingkungan teorema hanya untuk pernyataan yang memang memiliki
status matematis yang jelas. Bedakan definisi, lemma, teorema, dan bukti.
\end{frame}
\begin{frame}{Materi khusus}
LaTeX dapat menangani banyak bidang dengan paket yang tepat:
\begin{itemize}
\item ujian dengan kelas atau paket khusus;
\item formula kimia dengan \texttt{mhchem};
\item satuan dengan \texttt{siunitx};
\item diagram dan notasi bidang melalui paket yang didukung komunitas.
\end{itemize}
Selalu cek dokumentasi paket dan kebutuhan kompilasi sebelum mengadopsinya.
\end{frame}
\section{Pengerjaan dan finalisasi}
\begin{frame}{Integrasi seluruh komponen}
Urutan integrasi yang aman:
\begin{enumerate}
\item pastikan struktur bab dan subbab sudah stabil;
\item masukkan tabel, gambar, persamaan, dan rujukan silang;
\item aktifkan bibliografi, indeks, glossary, dan daftar isi;
\item kompilasi dari sumber utama dan simpan log yang bersih.
\end{enumerate}
\end{frame}
\begin{frame}{Pemeriksaan struktur}
\begin{itemize}
\item semua bagian penting memiliki heading yang sesuai;
\item urutan bab, lampiran, dan halaman awal mengikuti panduan;
\item setiap gambar dan tabel dirujuk dari teks;
\item tidak ada judul kosong, label ganda, atau berkas yang tertinggal.
\end{itemize}
\end{frame}
\begin{frame}{Pemeriksaan layout dan typography}
Cari dan perbaiki:
\begin{itemize}
\item teks melewati margin atau ruang kosong yang mengganggu;
\item judul terpisah dari paragraf atau caption dari objeknya;
\item font, spasi, tanda kutip, dan notasi berubah tanpa alasan;
\item halaman pembuka bab, header, footer, dan nomor halaman tidak konsisten.
\end{itemize}
\end{frame}
\begin{frame}[fragile]{Pemeriksaan referensi dan kompilasi}
\begin{lstlisting}[style=terminal]
latexmk -pdf main.tex
% periksa log untuk Undefined references dan citations
% kompilasi ulang sampai daftar isi dan nomor stabil
\end{lstlisting}
Perbaiki error pertama, bukan semua pesan sekaligus. Warnings tetap perlu
dibaca, terutama \texttt{overfull \textbackslash hbox}, referensi tak terdefinisi,
dan sitasi yang hilang.
\end{frame}
\begin{frame}{Konsultasi berdasarkan dokumen peserta}
Gunakan dokumen peserta sebagai bahan diagnosis:
\begin{itemize}
\item minta sumber minimal yang dapat dikompilasi;
\item pisahkan masalah isi, struktur, layout, dan toolchain;
\item ubah satu hal pada satu waktu lalu bandingkan PDF;
\item simpan solusi yang berulang sebagai perintah atau template.
\end{itemize}
\end{frame}
\begin{frame}{Checklist finalisasi}
\begin{enumerate}
\item PDF dapat dikompilasi dari sumber utama.
\item Daftar isi, daftar tabel, dan daftar gambar sudah diperbarui.
\item Semua sitasi, referensi silang, indeks, dan glossary terisi.
\item Tidak ada error, referensi hilang, atau teks keluar margin.
\item Folder proyek berisi berkas yang diperlukan dan instruksi kompilasi.
\end{enumerate}
\end{frame}
\begin{frame}{Latihan akhir}
Ambil salah satu dokumen peserta dan lakukan finalisasi singkat:
\begin{enumerate}
\item rapikan layout halaman dan gaya heading;
\item pisahkan isi menjadi beberapa berkas;
\item tambahkan satu referensi, satu rujukan silang, dan satu daftar otomatis;
\item jalankan kompilasi bersih dan catat perubahan.
\end{enumerate}
\end{frame}
\begin{frame}[plain]
\centering\Large Pertanyaan dan konsultasi\par\bigskip\normalsize
Simpan sumber, log, dan PDF final bersama-sama.
\end{frame}
\end{document}
